\section{Background and Threat Model}\label{sec:threat}

\subsection{Authorization in IEEE~2030.5 / CSIP}
% provenance: jrn_01KXT490ZD4444AWMRG5296GJA (threat model + CSIP function sets); lit passos2025tutorial.
IEEE~2030.5 uses mutual TLS with X.509 device certificates. A device's identity is the
certificate-derived LFDI/SFDI, and its authorization is scoped by a
\texttt{FunctionSetAssignments} resource binding it to one or more \texttt{DERProgram}s.
Physical control authority is carried by the \texttt{DERControlBase} modes: active-power capping
(\texttt{opModMaxLimW}), signed fixed active power (\texttt{opModFixedW}), reactive power
(\texttt{opModFixedVar}, \texttt{opModVoltVar}), volt-watt and frequency-watt curves, and the hard
connect and energize controls~\cite{passos2025tutorial}. The standard specifies certificate-based
authentication and function-set scoping; it does not mandate a single canonical access-control
object, and no accessible source establishes that it can bound the \emph{value} a control may
take. That distinction is where our design sits, and Section~\ref{sec:design} states it plainly
rather than implying native expressiveness.

\subsection{Six ways authority can be withdrawn, and what each reaches}
The literature tends to treat ``revocation'' as one thing. It is six, and they terminate different
layers, which is the whole reason a compromise can outlive the response to it:

\begin{enumerate}\itemsep2pt
\item \emph{PKI revocation} --- a CRL or OCSP status marks the certificate invalid.
\item \emph{Local identity denial} --- the server refuses the certificate-derived identity.
\item \emph{Authorization removal} --- \texttt{FunctionSetAssignments} or \texttt{DERProgram}
      membership changes.
\item \emph{Session termination} --- an established authenticated channel is closed or forced to
      revalidate.
\item \emph{Command cancellation} --- a scheduled or latched \texttt{DERControl} is withdrawn,
      superseded, or replaced by a safe default.
\item \emph{Identity replacement} --- the device is re-enrolled or re-keyed.
\end{enumerate}

Only the first is unavailable by construction. The motivating claim of this paper is therefore not
that IEEE~2030.5 offers no way to withdraw authority --- it offers the second, and CSIP builds
enrollment on it --- but that \textbf{the end-to-end time from deciding to contain a compromise to
the cessation of all effective physical authority is not bounded by any of them individually}.
Denying an identity refuses future authorization. On its own it neither closes an open session nor
retracts a control already issued, and Section~\ref{sec:results} measures how much of an attack
survives it.

\subsection{What the accessible standards material establishes}
% provenance: jrn_01KXT3JPAZ8RSGZK99X4G3HWPE (SAND2019-1490 4.1, primary-sourced, full text).
Certificate management differs across editions: IEEE~2030.5-2018~\cite{ieee20305-2018},
IEEE~2030.5-2023~\cite{ieee20305-2023} (which supersedes it, published December 2024), and the
profile in force~\cite{sunspec2018csip}. \textbf{We were unable to obtain IEEE~2030.5-2023.} It is
paywalled at roughly US\$255 and is not in the IEEE GET programme. Every claim below is evidenced
against the 2018 edition and the Sandia analysis of it~\cite{obert2019recommendations}, whose full
text is freely available; Table~\ref{tab:standards} gives each claim its source and its exposure,
and Section~S1 of the supplement reproduces the supporting quotations.

Three points are settled by that material and matter to the design.

\emph{Certificates do not expire, and revocation is prohibited for IEEE~2030.5 certificates.} The
Sandia analysis states that the PKI is distinguished by ``the use of non-expiring certificates and
the explicit prohibition of CRL and OCSP'', that ``once issued, a device certificate has an
indefinite lifetime'', and that ``if the device's private key is compromised, the device
certificate cannot be revoked''~\cite{obert2019recommendations}. The prohibition is specific:
OCSP ``may only be used to verify non-IEEE~2030.5 certificates''. We state the qualification
because an unqualified claim that the protocol forbids OCSP outright is wrong.

\emph{Local allow/deny listing is the sanctioned substitute, and is weaker than revocation.} It is
permitted, not required. The same analysis is explicit about its limits: ``a blacklist is
essentially a local CRL but without the process, rigor and verification behind its issuance'', and
``it is possible for a stolen device certificate that was disallowed in one region to be used in a
different region to gain access because it was only locally blacklisted''.

\emph{A short-lived operational credential is permitted, not prohibited.} Manufacturers following
the guidance ``are not required to limit the life of a device certificate''. Limited-life
certificates are in fact recommended, at a one-year renewal interval. CONTAINDER's contribution is
therefore not the short lifetime, which is neither novel nor disallowed, but the enforcement of
expiry at the session and command-effect layers, on which the accessible material is silent.

\begin{table}[t]
\centering
\caption{Standards claims the argument uses, the free source of record, and residual exposure.
SAND = Sandia SAND2019-1490~\cite{obert2019recommendations}, full text public; CSIP = Common
Smart Inverter Profile v2.1~\cite{sunspec2018csip}. \textbf{No row has been checked against
IEEE~2030.5-2023}, which could not be obtained.}
\label{tab:standards}
\footnotesize
\begin{tabular}{@{}p{0.04\columnwidth}p{0.36\columnwidth}p{0.24\columnwidth}p{0.22\columnwidth}@{}}
\toprule
ID & Claim & Free source & Residual risk \\
\midrule
S1 & Device certificate lifetime indefinite & SAND \S4.1, \S6.4, verbatim & 2023 clause unread \\
S2 & CRL prohibited for 2030.5 certificates & SAND \S4.1, verbatim & 2023 clause unread \\
S3 & OCSP prohibited \emph{for 2030.5 certificates}; permitted for others & SAND \S6.4, verbatim &
2023 clause unread \\
S4 & Local allow/deny listing permitted, not required, and unsynchronised & SAND \S4.1, \S4.1.2 &
low \\
S5 & Short-lived operational credential permitted & SAND \S6.4, verbatim & low \\
S6 & LFDI/SFDI derives from certificate & SAND \S3.1.1; CSIP \S5.2.1.2 & low; two sources \\
S7 & FunctionSetAssignments scopes \emph{functions}; numeric value bounds not established &
CSIP \S5.2.3.1 & \textbf{design depends on it}; see \S\ref{sec:design} \\
S8 & Scheduled DER controls persist & CSIP \S4.4.2, \S6.1.8.2 & \textbf{high}; 2023 drops
mandatory DERControl modes \\
\bottomrule
\end{tabular}
\end{table}

\subsection{Three separated layers}
We separate three layers the baseline conflates. \emph{L1, bootstrap identity}: the
hardware-rooted device certificate. \emph{L2, operational credential}: what authorizes live
sessions to an aggregator or DER management system. \emph{L3, authorization scope}: the
server-side binding from an identity to the function sets and values it may invoke. In baseline
IEEE~2030.5, L1 and L2 are the same indefinitely valid certificate, with no attestation gate. That
conflation, not any cryptographic weakness, is what makes a compromise persist.

\subsection{Compromise classes and what bounds each}
Retained authority differs by compromise class, and reporting one figure for all of them is the
error the pilot made. Table~\ref{tab:classes} gives the classes, the mechanism that bounds each,
and the resulting expectation. The derivations and their Monte-Carlo checks are in
\texttt{pkimodel/retention.py} and \texttt{tests/test\_retention.py}.

Two entries deserve emphasis because an earlier version of this paper got them wrong. For a
\emph{copied operational key} held by an adversary who is not resident on the device, attestation
has nothing to detect: renewal mints a fresh key whether or not anything is found, so retention is
bounded by the remaining epoch at $T/2$ in expectation, \emph{independently of the detection
probability}. For a \emph{persistent} compromise re-checked at each renewal with probability $p$,
the familiar $T/p$ holds only when the compromise arrives immediately after an issuance; under
uniform arrival within the epoch the expectation is $T(1/p - 1/2)$. At $T = 6$\,h and $p = 0.2$,
reporting $T/p$ for a copied key claims $30$\,h of retention where the true expectation is $3$\,h.

\begin{table}[t]
\centering
\caption{Compromise classes, the mechanism that bounds each, and expected retained authority.
$T$ is the credential lifetime and $p$ the per-renewal detection probability. Arrival is uniform
within the epoch unless stated.}
\label{tab:classes}
\footnotesize
\begin{tabular}{@{}p{0.30\columnwidth}p{0.34\columnwidth}p{0.24\columnwidth}@{}}
\toprule
Compromise class & Bounding mechanism & Expected retention \\
\midrule
Copied operational key & fresh key each epoch & $T/2$; \emph{independent of $p$} \\
Persistent, attested & attestation-gated renewal & $T(1/p - 1/2)$ \\
Persistent, unmeasured & none; scope only & unbounded \\
Stolen active session & session-age or per-command revalidation & $\min(T_{\text{age}},
T_{\text{reval}})$ \\
Persistent command & duration cap, cancellation, safe restore & $\min(T_{\text{dur}},
T_{\text{cancel}})$ \\
Stolen bootstrap key & hardware binding, enrollment policy & enrollment-controlled \\
Compromised gateway & gateway attestation, issuer policy & not bounded downstream \\
Compromised issuer & governance, HSM, audit & protection boundary \\
\bottomrule
\end{tabular}
\end{table}

\subsection{Assets, adversary, and goals}
\emph{Protected assets:} DER control integrity, fleet authorization, operational and attestation
keys, feeder stability and service continuity. \emph{Adversary capabilities:} steals an
operational key, compromises firmware, steals an active session, compromises a gateway or
aggregator, issues protocol-conformant commands within its authorization, and observes authorized
telemetry. It does not break TLS or the underlying cryptography. \emph{Exclusions:} breaking
standard cryptography, physical destruction, and compromising both the root issuer and all audit
controls, which we treat as the protection boundary rather than a defended case.
\emph{Security goals:} least authority, bounded new-session authority, bounded active-session
authority, bounded command-effect duration, fresh-key containment, replay resistance, safe
degradation, auditability. \emph{Non-goals:} defending against malicious root issuance without
added governance, detecting every hardware attack, replacing grid protection, and proving safety
for every feeder. Attestation detects only what the evidence measures and the reference-value
policy distinguishes; a stolen operational key on a device with intact firmware passes it, which
is why fresh-key binding rather than attestation is what bounds that class.
